% q2.tex
\begin{frame}{q2 問題}
  \begin{alertblock}{問い}
    $A \le_{\mathrm{m}} B$ であり、かつ $B$ が正規言語であるとする。
    このとき、$A$ は必ず正規言語であると言えるか？
  \end{alertblock}
  \begin{block}{答え}
    $A$ は必ず正規言語であるとは限らない。
  \end{block}
\end{frame}


\begin{frame}{具体例の構成}
  \begin{itemize}
    \item $B$ を正規言語とする（例：$B = \{0^*1^*\}$）。
    \item $A$ を非正規言語とする（例：$A = \{0^n1^n \mid n \ge 0\}$）。
    \item $A$ から $B$ への写像帰着を構築することが可能である。
  \end{itemize}
\end{frame}

\begin{frame}{写像帰着の構成}
    $A$ から $B$ への写像帰着を構築するための関数 $f$ を以下のように定義する。

    \begin{exampleblock}{関数 $f$ の定義}
        入力 $x$ に対して：
        \begin{enumerate}
            \item $x$ が $0^n1^n$ の形式であれば、\textbf{01 を出力する}。
            \item そうでなければ、\textbf{10 を出力する}。
        \end{enumerate}
    \end{exampleblock}

    \vfill
    \begin{itemize}
        \item この関数は計算可能である。
        \item $f(x) = 01$ は $B$ に属し、$f(x) = 10$ は $B$ に属さないため、帰着の条件 $x \in A \iff f(x) \in B$ を満たす。
    \end{itemize}
\end{frame}

\begin{frame}{結論}
    以上の構成により、$A$ は $B$ へ写像帰着可能であるが、$A$ は非正規言語であるため、$A$ は必ずしも正規言語であるとは限らないことが示された。
\end{frame}
